\documentclass[11pt,a4paper]{report}

\usepackage[T1]{fontenc}
\usepackage[utf8]{inputenc}
\usepackage[french]{babel}
\usepackage{geometry}
\usepackage{graphicx}
\usepackage{hyperref}
\usepackage{booktabs}
\usepackage{longtable}
\usepackage{array}
\usepackage{xcolor}
\usepackage{enumitem}
\usepackage{fancyhdr}

\geometry{margin=2.5cm}
\hypersetup{
  colorlinks=true,
  linkcolor=blue!60!black,
  urlcolor=blue!60!black,
  pdftitle={Rapport détaillé — Site Target Agency},
  pdfauthor={Target Agency / Wetarget}
}

\pagestyle{fancy}
\fancyhf{}
\fancyhead[L]{\small Rapport Target Agency}
\fancyhead[R]{\small\today}
\fancyfoot[C]{\thepage}

\title{
  \vspace{1cm}
  {\Huge\bfseries Rapport détaillé du site}\\[0.4cm]
  {\Large Target Agency — Projet \texttt{wetarget}}\\[0.8cm]
  {\large État des lieux complet à date du 5 septembre 2026}
}
\author{Document généré à partir de l'analyse du dépôt de code}
\date{}

\begin{document}

\maketitle
\thispagestyle{empty}
\newpage

\tableofcontents
\newpage

%==============================================================================
\chapter{Résumé exécutif}
%==============================================================================

Le projet \textbf{wetarget} est une application \textbf{Next.js 16.3.4} hébergeant le site vitrine de \textbf{Target Agency}, agence de marketing digital et de communication. Le site public provient d'un \textbf{export Framer} (NoCodeXport) servi tel quel via une route catch-all. Un \textbf{panneau d'administration} custom (shadcn/ui + Tailwind CSS 4) a été ajouté pour piloter l'agence. Une base \textbf{Supabase} stocke les messages du formulaire de contact.

\begin{itemize}[leftmargin=*]
  \item \textbf{51 pages HTML} exportées dans \texttt{content/}
  \item \textbf{410 fichiers} d'assets statiques dans \texttt{public/assets/}
  \item \textbf{5 dashboards} par activité + \textbf{1 vue globale} agence
  \item \textbf{1 table Supabase} : \texttt{contact\_submissions}
  \item Stack : React 19, TypeScript 5, Recharts, Radix UI, Jose (JWT admin)
\end{itemize}

%==============================================================================
\chapter{Architecture technique}
%==============================================================================

\section{Stack logicielle}

\begin{table}[h]
\centering
\begin{tabular}{@{}ll@{}}
\toprule
\textbf{Composant} & \textbf{Version / détail} \\
\midrule
Next.js & 16.3.4 (Turbopack en dev) \\
React & 19.2.8 \\
TypeScript & 5.x \\
Tailwind CSS & 4.3.3 (scopé à l'admin) \\
Supabase JS & 2.115.0 + @supabase/ssr \\
Authentification admin & JWT (jose), cookie HTTP-only \\
Graphiques admin & Recharts 3.10.1 \\
UI admin & shadcn/ui (Radix + CVA) \\
\bottomrule
\end{tabular}
\caption{Stack principale du projet}
\end{table}

\section{Structure des dossiers}

\begin{description}[style=nextline, leftmargin=3cm]
  \item[\texttt{content/}] Pages HTML exportées Framer (site public)
  \item[\texttt{public/assets/}] Images, polices, vidéos, scripts et CMS Framer
  \item[\texttt{app/[[...slug]]/}] Route catch-all servant le HTML Framer
  \item[\texttt{app/admin/}] Panneau d'administration (login + dashboards)
  \item[\texttt{app/api/framercms/}] API pour requêtes range sur fichiers \texttt{.framercms}
  \item[\texttt{components/}] Composants React (UI shadcn + admin)
  \item[\texttt{lib/}] Utilitaires, Supabase, auth admin, config dashboards
  \item[\texttt{supabase/}] Migration SQL et configuration CLI
  \item[\texttt{types/}] Types TypeScript générés (base de données)
\end{description}

\section{Routage et rendu du site public}

Le fichier \texttt{app/[[...slug]]/route.ts} :
\begin{itemize}[leftmargin=*]
  \item Résout chaque URL vers un fichier \texttt{content/.../index.html}
  \item Génère les paramètres statiques via \texttt{generateStaticParams}
  \item Retourne le HTML brut avec \texttt{Content-Type: text/html}
  \item En cas d'absence de page : sert \texttt{content/404.html} (si présent)
\end{itemize}

\textbf{Note :} il n'existe pas de \texttt{content/index.html} à la racine. L'accueil principal en français est accessible via \texttt{/fr}. La route \texttt{/} dépend de la résolution Framer / redirection côté export.

\section{Configuration Next.js (\texttt{next.config.ts})}

\begin{itemize}[leftmargin=*]
  \item \textbf{output standalone} : activé hors Vercel (désactivé sur Vercel pour éviter une erreur NFT)
  \item \textbf{images unoptimized} : true (compatibilité export statique)
  \item \textbf{Rewrite} : \texttt{/assets/animate/*?range=} $\rightarrow$ API framercms (CMS Framer)
  \item \textbf{outputFileTracingIncludes} : inclut \texttt{content/**} et fichiers \texttt{.framercms}
\end{itemize}

\section{Middleware}

Le fichier \texttt{middleware.ts} délègue à \texttt{lib/supabase/middleware.ts} qui :
\begin{enumerate}[leftmargin=*]
  \item Protège \texttt{/admin/*} (sauf \texttt{/admin/login}) via session JWT admin
  \item Redirige les utilisateurs connectés depuis la page login vers \texttt{/admin}
  \item Rafraîchit la session Supabase Auth (cookies) sur les autres routes
\end{enumerate}

%==============================================================================
\chapter{Site public — Target Agency}
%==============================================================================

\section{Identité et positionnement}

\begin{itemize}[leftmargin=*]
  \item \textbf{Nom :} Target Agency
  \item \textbf{Titre SEO :} «~Agence de Marketing digital \& de communication~»
  \item \textbf{Slogan hero (FR) :} «~Votre image. Notre obsession.~»
  \item \textbf{Promesse :} Transformer les ambitions en identités de marque fortes, créant confiance et empreinte durable
  \item \textbf{Fondateur mentionné :} Adam Fisli
  \item \textbf{Contact affiché :} +33 6 09 88 46 60
\end{itemize}

\section{Les cinq activités / services}

\begin{table}[h]
\centering
\small
\begin{tabular}{@{}p{4.5cm}p{9cm}@{}}
\toprule
\textbf{Service} & \textbf{Contenu principal} \\
\midrule
Branding \& Identité visuelle &
Logo, charte graphique, identité de marque, supports de communication \\

Création de sites web &
UX/UI, design responsive, optimisation UX, mise en ligne et suivi \\

Stratégie digitale \& Réseaux sociaux &
Stratégie de com, gestion des réseaux, planification, analytics \\

Acquisition \& Publicité &
Campagnes ads, ciblage, optimisation performances, suivi ROAS \\

Automatisation \& Intelligence artificielle &
Automatisation processus, connexion d'outils, workflows, optimisation ops \\
\bottomrule
\end{tabular}
\caption{Offre de services telle qu'affichée sur la landing page}
\end{table}

\section{Forces mises en avant sur le site}

\begin{enumerate}[leftmargin=*]
  \item Expertise 360° — un seul partenaire pour toute la présence digitale
  \item Projets livrés clés en main (pas seulement des maquettes)
  \item De la stratégie à l'exécution (pilotage au quotidien)
  \item Résultats concrets : visibilité, engagement, conversions
  \item Collaboration transparente avec jalons et visibilité
\end{enumerate}

\section{Processus en 4 étapes}

\begin{enumerate}[leftmargin=*]
  \item Découverte \& Analyse
  \item Stratégie \& Créativité
  \item Création \& Développement
  \item Lancement \& Croissance
\end{enumerate}

\section{Portfolio — études de cas}

\begin{table}[h]
\centering
\begin{tabular}{@{}lll@{}}
\toprule
\textbf{Client / projet} & \textbf{Route EN} & \textbf{Route FR} \\
\midrule
Media \& Brokerage By Reign & \texttt{/cases/mb-byreign} & \texttt{/fr/cases/mb-byreign} \\
Luxkey & \texttt{/cases/luxkey} & \texttt{/fr/cases/luxkey} \\
Latifa B. & \texttt{/cases/latifa-b} & \texttt{/fr/cases/latifa-b} \\
Les Gourmandises de Nadj & \texttt{/cases/les-gourmandises-de-nadj} & \texttt{/fr/cases/les-gourmandises-de-nadj} \\
L'Héritage Français & \texttt{/cases/lheritage-francais} & \texttt{/fr/cases/lheritage-francais} \\
Liga Del Pueblo & \texttt{/cases/liga-del-pueblo} & \texttt{/fr/cases/liga-del-pueblo} \\
\bottomrule
\end{tabular}
\caption{Études de cas disponibles (versions EN et FR)}
\end{table}

\section{Blog — articles (12 sujets × 2 langues)}

Thématique dominante : \textbf{systèmes d'IA}, automatisation et workflows en production (contenu hérité du template Framer, orienté tech/IA en anglais).

\begin{longtable}{@{}p{10.5cm}@{}}
\toprule
\textbf{Slug de l'article} \\
\midrule
\endhead
\texttt{beyond-the-text-box-engineering-autonomous-workflows-over-agency-free-chat} \\
\texttt{building-decision-engines-instead-of-chat-interfaces} \\
\texttt{designing-multi-layer-ai-systems-for-real-world-operations} \\
\texttt{from-conversation-to-conclusion-why-your-product-needs-outcomes-not-dialogue} \\
\texttt{from-fragmented-data-to-unified-intelligence-engines} \\
\texttt{from-manual-bottlenecks-to-autonomous-growth-at-scale} \\
\texttt{from-prompts-to-pipelines-architecting-intentional-intelligence-over-reactive-replies} \\
\texttt{from-workflow-chaos-to-structured-automation-pipelines} \\
\texttt{stop-prompting-start-predicting-transitioning-from-llms-to-logic-layers} \\
\texttt{the-autonomy-leap-moving-from-conversation-as-interface-to-outcome-as-service} \\
\texttt{why-most-ai-deployments-fail-in-production-environments} \\
\bottomrule
\end{longtable}

Routes : \texttt{/blog/<slug>} et \texttt{/fr/blog/<slug>}

%==============================================================================
\chapter{Inventaire complet des pages (51 routes)}
%==============================================================================

\section{Pages principales (anglais)}

\begin{itemize}[leftmargin=*]
  \item \texttt{/about-us} — À propos
  \item \texttt{/contacts} — Contact
  \item \texttt{/career} — Carrières
  \item \texttt{/blog} — Index blog
  \item \texttt{/cases} — Index portfolio
  \item \texttt{/legal} — Mentions légales
  \item \texttt{/policy} — Politique de confidentialité
  \item \texttt{/terms} — Conditions générales
\end{itemize}

\section{Pages principales (français — préfixe \texttt{/fr})}

\begin{itemize}[leftmargin=*]
  \item \texttt{/fr} — \textbf{Accueil principal (landing page FR)}
  \item \texttt{/fr/about-us}
  \item \texttt{/fr/contacts}
  \item \texttt{/fr/career}
  \item \texttt{/fr/blog}
  \item \texttt{/fr/cases}
  \item \texttt{/fr/legal}
  \item \texttt{/fr/policy}
  \item \texttt{/fr/terms}
\end{itemize}

\section{Total}

\begin{center}
\textbf{51 fichiers \texttt{index.html}} répartis en pages EN, FR, blog et cas clients.
\end{center}

%==============================================================================
\chapter{Assets et export Framer}
%==============================================================================

\section{Contenu statique}

\begin{itemize}[leftmargin=*]
  \item \textbf{410 fichiers} dans \texttt{public/assets/}
  \item Polices : Inter, Inter Display, Plus Jakarta Sans, JetBrains Mono
  \item Scripts Framer : modules \texttt{.mjs} dans \texttt{public/assets/animate/}
  \item Fichiers CMS : \texttt{.framercms} pour contenu dynamique Framer
  \item Métadonnées SEO : \texttt{public/assets/meta/*.json}
  \item Images, icônes, vidéos exportées
\end{itemize}

\section{Particularité technique Framer CMS}

Les requêtes \texttt{GET /assets/animate/*.framercms?range=...} sont réécrites vers \texttt{/api/framercms/[...file]} pour renvoyer uniquement la plage d'octets demandée (compatibilité Vercel / hébergement).

%==============================================================================
\chapter{Base de données Supabase}
%==============================================================================

\section{Schéma (\texttt{contact\_submissions})}

\begin{table}[h]
\centering
\small
\begin{tabular}{@{}llp{6cm}@{}}
\toprule
\textbf{Colonne} & \textbf{Type} & \textbf{Description} \\
\midrule
\texttt{id} & uuid & Clé primaire (gen\_random\_uuid) \\
\texttt{full\_name} & text & Nom complet (requis) \\
\texttt{email} & text & Email (requis) \\
\texttt{phone} & text & Téléphone (optionnel) \\
\texttt{company} & text & Entreprise (optionnel) \\
\texttt{message} & text & Message (requis) \\
\texttt{source\_page} & text & Page d'origine du formulaire \\
\texttt{locale} & text & Langue (défaut : \texttt{fr}) \\
\texttt{status} & text & \texttt{new}, \texttt{read}, \texttt{replied}, \texttt{archived} \\
\texttt{created\_at} & timestamptz & Date de création \\
\texttt{updated\_at} & timestamptz & Mise à jour auto (trigger) \\
\bottomrule
\end{tabular}
\end{table}

\section{Sécurité (RLS)}

\begin{itemize}[leftmargin=*]
  \item \textbf{INSERT} : autorisé pour \texttt{anon} et \texttt{authenticated} (formulaire public)
  \item \textbf{SELECT / UPDATE} : réservé aux utilisateurs \texttt{authenticated}
  \item L'admin actuel utilise la clé secrète Supabase côté serveur
\end{itemize}

\section{Clients Supabase dans le code}

\begin{itemize}[leftmargin=*]
  \item \texttt{lib/supabase/client.ts} — navigateur
  \item \texttt{lib/supabase/server.ts} — Server Components
  \item \texttt{lib/supabase/admin.ts} — opérations admin (clé secrète \texttt{sb\_secret\_*})
  \item \texttt{lib/supabase/health.ts} — test de connexion + comptage contacts
  \item \texttt{lib/supabase/contacts.ts} — lecture messages + stats conversion
\end{itemize}

\section{Variables d'environnement}

\begin{table}[h]
\centering
\small
\begin{tabular}{@{}lp{8cm}@{}}
\toprule
\textbf{Variable} & \textbf{Rôle} \\
\midrule
\texttt{NEXT\_PUBLIC\_SUPABASE\_URL} & URL du projet Supabase \\
\texttt{NEXT\_PUBLIC\_SUPABASE\_PUBLISHABLE\_KEY} & Clé publique (anon) \\
\texttt{SUPABASE\_SECRET\_KEY} & Clé secrète serveur (jamais exposée) \\
\texttt{ADMIN\_USERS} & Liste \texttt{user:pass} séparés par virgule \\
\texttt{ADMIN\_SESSION\_SECRET} & Secret JWT admin ($\geq$ 32 caractères) \\
\bottomrule
\end{tabular}
\end{table}

%==============================================================================
\chapter{Panneau d'administration (\texttt{/admin})}
%==============================================================================

\section{Authentification}

\begin{itemize}[leftmargin=*]
  \item Route login : \texttt{/admin/login}
  \item Utilisateurs configurés : \texttt{adam} et \texttt{mahdi} (mot de passe via \texttt{.env.local})
  \item Session : cookie HTTP-only \texttt{admin\_session}, JWT signé HS256 (jose)
  \item Durée : définie dans \texttt{lib/admin/constants.ts}
  \item Protection : middleware sur toutes les routes \texttt{/admin/*} sauf login
\end{itemize}

\section{Interface (design shadcn/ui)}

\begin{itemize}[leftmargin=*]
  \item Styles : \texttt{app/admin/admin.css} (Tailwind 4, variables CSS shadcn)
  \item Layout : sidebar + header (recherche, thème) + zone principale
  \item Composants KPI : cartes avec sparkline corail, bouton Détails, pourcentage vert/rouge
  \item Graphiques : Recharts (area, bar, line, donut/pie)
  \item Tableaux : paiements mock, messages contact Supabase, équipe
\end{itemize}

\section{Vue globale — \texttt{/admin}}

\begin{table}[h]
\centering
\begin{tabular}{@{}lp{9cm}@{}}
\toprule
\textbf{Élément} & \textbf{Détail} \\
\midrule
CA total du mois & 108\,250\,€ (données mock) \\
Projets en cours & 24 (mock) \\
Leads contact & Données \textbf{réelles} Supabase \\
Conversion lead $\rightarrow$ client & Calculée : (répondu + archivé) / total \\
Graphique donut & Branding 30\%, Web 25\%, Social 20\%, Pub 15\%, IA 10\% \\
Accès rapides & Liens vers les 5 dashboards activité \\
\bottomrule
\end{tabular}
\caption{Vue globale agence}
\end{table}

\section{Dashboards par activité}

\begin{longtable}{@{}llp{7cm}@{}}
\toprule
\textbf{ID} & \textbf{Label} & \textbf{KPI spécifiques (exemples)} \\
\midrule
\endhead
\texttt{branding} & Branding & Projets actifs, identités livrées, satisfaction \\
\texttt{sites-web} & Sites web & Sites en prod, pages déployées, score perf. \\
\texttt{reseaux-sociaux} & Réseaux sociaux & Comptes gérés, publications, engagement \\
\texttt{publicite} & Publicité & Campagnes actives, leads, ROAS \\
\texttt{automatisation} & Automatisation \& IA & Workflows actifs, heures économisées, outils connectés \\
\bottomrule
\caption{Dashboards \texttt{/admin/dashboard/[id]}}
\end{longtable}

Chaque dashboard contient :
\begin{itemize}[leftmargin=*]
  \item 3 cartes KPI + 1 carte revenu avec mini graphique
  \item Graphique area chart (activité mensuelle)
  \item Graphique bar chart
  \item Table facturation (données mock) + messages contact Supabase
  \item Liste membres équipe (Adam Fisli, Mahdi Benali, Sophie Leroy)
  \item Onglets : Vue d'ensemble, Analytics, Rapports, Notifications
\end{itemize}

\section{Composants React admin}

\begin{itemize}[leftmargin=*]
  \item \texttt{components/admin/admin-shell.tsx} — Sidebar + header
  \item \texttt{components/admin/agency-overview.tsx} — Vue globale
  \item \texttt{components/admin/dashboard-overview.tsx} — Dashboard activité
  \item \texttt{components/admin/kpi-card.tsx} — Carte KPI style SaaS
  \item \texttt{components/ui/*} — button, card, tabs, input, badge, avatar, etc.
  \item \texttt{components/ui/hero-195.tsx} — Composant hero marketing (non affiché sur dashboards)
\end{itemize}

%==============================================================================
\chapter{API interne}
%==============================================================================

\section{\texttt{GET /api/framercms/[...file]}}

Route handler pour servir les fichiers \texttt{.framercms} avec support du header/query \texttt{range} (requis par le runtime Framer pour le CMS).

%==============================================================================
\chapter{Scripts npm}
%==============================================================================

\begin{table}[h]
\centering
\begin{tabular}{@{}lp{9cm}@{}}
\toprule
\textbf{Commande} & \textbf{Action} \\
\midrule
\texttt{npm run dev} & Serveur de développement (port 3000) \\
\texttt{npm run build} & Build production (SSG + routes admin dynamiques) \\
\texttt{npm run start} & Serveur production \\
\texttt{npm run lint} & ESLint \\
\texttt{npm run supabase:types} & Génère \texttt{types/database.ts} \\
\bottomrule
\end{tabular}
\end{table}

%==============================================================================
\chapter{Limites actuelles et évolutions possibles}
%==============================================================================

\section{Données mock vs réelles}

\begin{itemize}[leftmargin=*]
  \item CA, projets, campagnes, paiements : \textbf{données fictives} dans l'admin
  \item Seuls les \textbf{messages contact Supabase} sont branchés en live
  \item Le formulaire Framer n'est pas encore connecté à Supabase côté code Next (insertion à implémenter)
\end{itemize}

\section{Contenu hérité du template}

\begin{itemize}[leftmargin=*]
  \item Page \texttt{/about-us} : contenu orienté «~AI Studio~» en anglais (template Framer)
  \item Blog : articles IA/automation en anglais, peu alignés avec l'offre FR marketing
  \item Numéros de téléphone placeholder (+1 555...) présents dans certaines métadonnées
\end{itemize}

\section{Fonctionnalités admin à venir (suggérées)}

\begin{enumerate}[leftmargin=*]
  \item Inbox messages contact avec changement de statut
  \item Pipeline projets / clients par activité
  \item Connexion CRM ou facturation réelle
  \item Calendrier éditorial réseaux sociaux
  \item Monitoring uptime des sites clients
  \item Brancher le formulaire contact Framer $\rightarrow$ Supabase
\end{enumerate}

\section{Points techniques}

\begin{itemize}[leftmargin=*]
  \item Pas de \texttt{content/index.html} : accueil EN à \texttt{/} potentiellement absent ; accueil FR sur \texttt{/fr}
  \item Middleware Next.js 16 : convention \texttt{middleware} dépréciée (migration \texttt{proxy} recommandée)
  \item Admin users : mots de passe en clair dans \texttt{.env.local} (acceptable en local, à remplacer en prod)
\end{itemize}

%==============================================================================
\chapter{Annexe — Arborescence des routes Next.js (App Router)}
%==============================================================================

\begin{verbatim}
app/
  layout.tsx                    # Layout racine (html lang=fr)
  [[...slug]]/route.ts          # Site Framer (toutes pages publiques)
  api/framercms/[...file]/route.ts
  admin/
    layout.tsx                  # Import admin.css
    login/
      page.tsx                  # Page connexion
      login-form.tsx            # Formulaire shadcn
      actions.ts                # Server action login/logout
    (protected)/
      layout.tsx                # AdminShell + auth check
      page.tsx                  # Vue globale agence
      dashboard/[dashboardId]/
        page.tsx                # Dashboard par activite
\end{verbatim}

\vfill
\begin{center}
\small\textit{Fin du rapport — Projet wetarget / Target Agency}
\end{center}

\end{document}
